% Draft Threats-to-Validity insert for external variability study (RQ-EXT-E).
% Merge manually into paper/threats_to_validity.tex after Phase 2 if study is run.
% Replaces or supplements the single-sentence external feasibility probe.

\paragraph{External inter-execution variability (exploratory).}
\label{sec:threats:external-variability}

A separate exploratory study (HumanEval+; protocol \texttt{EXTERNAL\_VARIABILITY\_PROTOCOL.md}) addresses whether functionally valid \emph{generated} solutions on an independent benchmark exhibit output instability under repeated execution, without INVERT instrumentation contracts.
It employs a frozen output-stability detector that is distinct from the Class~E \texttt{deterministic\_randomized} auditor and does not use \texttt{ItemProcessor} or \texttt{visit\_fn}.
Prior feasibility probes showed that the frozen Class~E detector abstains on arbitrary Python; the external study therefore does not transfer Class~E pole-recovery claims.

The study asks descriptive prevalence questions (stable vs.\ variable vs.\ flaky among EvalPlus-valid completions) and whether one-shot functional pass misses repeated-run instability.
Canonical reference solutions are included as a sanity baseline only.
Outcomes, if reported, appear in the appendix and do not enter confirmatory Core~v2 recovery tables.

This evidence can support bounded statements about reproducibility risk on a standard code-generation benchmark.
It does \emph{not} validate Classes~C or~D externally, establish interface-agnostic recovery for quantity or order, or imply that INVERT's synthetic trace-contract benchmarks generalize to arbitrary repositories.
Adapter-based wrapping of external code into INVERT harnesses remains out of scope because it would reintroduce co-design.

% Optional: retain existing sentence on trace-contract probes if study not yet run:
% A feasibility probe documented trace-contract boundaries (\externalclosure{} in the replication package).
